\slajdid{02-s002}
\begin{frame}[label=02-s002]
\gusttekst
Kombinacione mreže su digitalni sistemi kod kojih izlazi zavise samo od stanja ulaza. ???

Od kojih stanja ulaza? Jučerašnjih? Od onih koja su bila pre 100ms? Bivših, sadašnjih, budućih?
Moramo dodati vremensku odrednicu!

Kombinacione mreže su digitalni sistemi kod kojih izlazi zavise samo od \alert{trenutnog} stanja ulaza.

Trenutnog? Pojam koji često koristimo ne ulazeći dublje u njegovo značenje.
Ako su se ulazi promenili u nekom trenutku vremena da li su se i izlazi \alert{trenutno} promenili.

Ako kažemo da se nešto trenutno dešava obično podrazumevamo neke vremenske margine pošto u realnom fizičkom svetu ne postoji pojava, promena, koja se zaista dešava trenutno, odnosno u beskonačno kratkom vremenskom intervalu.

Šta može da se desi i šta se dešava u digitalnom sistemu realizovanom kao kombinaciona mreža?

\alert{Ako se u nekom trenutku promene ulazi neće se odmah promeniti i izlazi. Postoji vremenski interval u kojem gornja tvrdnja ne važi. Stanja izlaza nam nisu poznata u tom vremenskom intervalu.}

Zašto?

Logička kola ne mogu trenutno da „reaguju“ na promenu ulaznih promenljivih. Fizika. Izlazi će se promeniti tek posle nekog vremena. \alert{Kašnjenje logičkog kola.}
\end{frame}
